\begin{corollary}
	Если $g_1, g_2 \in R[a; b]$, то $g_1 \cdot g_2 \in R[a; b]$
\end{corollary}

\begin{proof}
	Положим $f(x) := x^2$. Тогда $f$ - интегрируема по Риману на всем $\R$. Остаётся заметить следующее равенство:
	\[
		g_1 \cdot g_2 = \frac{1}{4}((g_1 + g_2)^2 - (g_1 - g_2)^2)
	\]
\end{proof}

\begin{corollary}
	Если $g \in R[a; b]$, то $|g| \in R[a; b]$, причём
	\[
		\left|\int_a^b g(x)dx\right| \le \int_a^b |g(x)|dx
	\]
\end{corollary}

\begin{proof}
	Положим $f(x) := |x|$, $c := \sgn \int_a^b g(x)dx$. Тогда
	\[
		\left|\int_a^b g(x)dx\right| = c\int_a^b g(x)dx = \int_a^b cg(x)dx \le \int_a^b |g(x)|dx
	\]
\end{proof}

\begin{note}
	Все полученные результаты про свойства справедливы и для интеграла Римана-Стилтьеса.
\end{note}

\begin{note}
	Можно также доказать, что из $f \in R(\alpha_1, [a; b])$ и $f \in R(\alpha_2, [a; b])$ следует, что $f \in R(\alpha_1 + \alpha_2, [a; b])$, причём
	\[
		\int_a^b f(x) d((\alpha_1 + \alpha_2)x) = \int_a^b f(x) d(\alpha_1(x)) + \int_a^b f(x)d(\alpha_2(x))
	\]
	Поставив <<$-$>> вместо <<$+$>>, получим определение интеграла Римана-Стилтьеcа для $\alpha = \alpha_1 - \alpha_2$, то есть функции ограниченной вариации. Можно показать, что корректность определения и все свойства будут работать и т.д. и т.п.
\end{note}

\subsection{Другие определения интегралов}

\begin{definition}
	Пусть $P:= (a=x_0<x_1<...<x_n=b)\text{-упорядоченное множество точек},\ \{t_i\}_{i=1}^n:\ x_{i-1} \leq t_i\leq x_i\ \forall i=1,..., n$ . \textit{Интегральной суммой} называется 
	\[
		S(P, f, \{t_i\})=\suml_{i = 1}^n f(t_i) \Delta x_i
	\]
\end{definition}

\begin{note}
	То есть у нас опять сумма площадей прямоугольников, но в этот раз они не вписаны, и не описаны около данной функции, а где-то между.
\end{note}

\begin{definition}
	Число $I$ называется \textit{пределом интегральных сумм}, $I = \liml_{\Delta(P) \to 0}S(P, f, \{t_i\})$, если
	\[
		(\forall \eps > 0)\ (\exists \delta > 0): \ (\forall P, \Delta(P) < \delta\ ) \ ( \forall \{t_i\},\ t_i \in [x_{i - 1}; x_i]\  )\hookrightarrow |S(P, f, \{t_i\}) - I| < \eps
	\]
\end{definition}

\begin{lemma}
	Если $\exists \liml_{\Delta(P) \to 0} S(P, f, \{t_i\})$, то $f$ ограничена на $[a; b]$.
\end{lemma}

\begin{proof}
	От противного. Не умаляя общности, пусть $f$ неограничена сверху на $[a; b]$. Покажем, что предела интегральных сумм не найдётся. Рассмотрим любое разбиение $P$, для которого $\Delta(P) < \frac{1}{n}$. В силу неограниченности\footnote{Напомним, что $M_k$ - супремум значений функции на $k$-ом отрезке разбиения}
	\[
		\exists k \such M_k = +\infty \Ra \exists t_k \in [x_{k - 1}; x_k]\ \ f(t_k) > \frac{n + |S'(P, f, \{t_i\})|}{\Delta x_k}
	\]
	Где  $S'(P, f, \{t_i\})$ это интегральная сумма всех отрезков, кроме k-ого. То есть по определению $S$:
	\[
		S(P, f, \{t_i\}) = S'(P, f, \{t_i\}) + f(t_k) \Delta x_k
	\]
	Оценим снизу модуль интегральной суммы $S$:
	\[
		|S(P, f, \{t_i\})| \ge |f(t_k) \Delta x_k| - |S'(P, f, \{t_i\})| = f(t_k)\Delta x_k - |S'(P, f, \{t_i\})| > n
	\]
	Таким образом модуль суммы построенного разбиения, с диаметром $\frac{1}{n}$ больше  $n$, а значит мы нашли построили разбиения, у которых при устремлении диаметра к нулю, модуль интегральной суммы неограничен. Значит предела нет, противоречие
\end{proof}

\begin{theorem} (Интеграл как предел интегральных сумм)
	$f \in R[a; b]$ тогда и только тогда, когда существует $\liml_{\Delta(P) \to 0} S(P, f, \{t_i\})$, причём
	\[
		\int_a^b f(x)dx = \liml_{\Delta(P) \to 0} S(P, f, \{t_i\})
	\]
\end{theorem}

\begin{proof}~
	\begin{itemize}
		\item $\La$ (Достаточность). По условию  $\exists I = \liml_{ \Delta(P) \to 0} S(P, f, \{t_i\})$ или же:
		\[
			(\forall \eps > 0)\ (\exists \delta > 0) : (\forall P, \Delta(P) < \delta)\ (\ \forall \{t_i\},\ t_i \in [x_{i - 1}; x_i])\hookrightarrow |S(P, f, \{t_i\}) - I| < {\eps}
		\]
		Тогда, если зафиксировать $P$:
		\[
			I-\eps < \suml_{i=  1}^nf(t_i)\Delta x_i < I+\eps
		\]
		Подставив в неравенство инфинумы и супремумы значений на отрезках, получим посередине суммы Дабру, а знак, в силу предельного перехода заменится на нестрогий:
		\[
			I-\eps \leq L(P, f) \leq U(P, f) \leq I+\eps
		\]
		То есть по критерию интегрируемости по Риману $f$ - интегрируема, т.к. разность верхней и нижней сумм $\leq 2\eps$, а так как интеграл Римана всегда лежит между суммами Дабру, неравенство гарантирует, что он и есть $I$.
		\item $\Ra$ (Необходимость) $f \in R[a,b] \Ra(\forall \eps > 0)(\exists P^*):$
		\[
		\System{
			&\ U(P^*, f) < \int_a^bf(x)dx + \frac{\eps}{4}
			\\
			&L(P^*, f) > \int_a^bf(x)dx - \frac{\eps}{4}
		}
		\]
		Дополнительно предположим, что $f(x) \geq0$, а также введём обозначения:
		\begin{itemize}
			\item$
			M=\sup_{x\in [a,b]} f(x)$
			\item ${P*:=x_0^*<x_1^*<...<x_{n^*}^*=b}$ 
			\item $\delta_1 := \frac{\eps}{8\cdot M\cdot n^*} > 0$
		\end{itemize}
	
		То есть $n^*$ - количество отрезков разбиения $P^*$. \\
		Теперь, мы можем изучить устройство $\forall P, \Delta P < \delta_1$. Разобьём сумму разбиения на две части, в первой - отрезки содержащие точки из $P^*$(кроме концов отрезка), а во второй - остальные. Тогда первую сумму можно оценить сверху так: каждое слагаемое из этой суммы - отрезок содержащий $x_i^*$, умноженный на свой супремум, количество таких отрезков оценивается $2(n^*-1)$, т.к. каждое $x_i^*$ не более чем в двух, а супремум и длина по определению оцениваются через $M$ и $\delta_1$. Со второй суммой всё ещё проще. Раз все отрезки из неё не содержат точек разбиения $P^*$, то каждый из них содержится в каком-то отрезке разбиения $P*$. Заметим, что если заменить все отрезки разбиения $P$ содержащиеся в данном отрезке из $P^*$ заменить на него, то сумма не уменьшится (т.к. супремум не уменьшится, и суммарная длина тоже). Значит все отрезки второй суммы, можно буквально заменить на совокупность всех отрезков из $P^*$, то есть просто верхнюю сумму $P^*$. 
		\begin{multline*}
			(\forall P, \Delta (P) < \delta_1) \quad U(P, f) = \suml_{\substack{k: \exists  x_i^*\in[x_{k-1}, x_k] \\ i \neq0,n}} M_i\Delta x_i + \suml_{\substack{k: \nexists  x_i^*\in[x_{k-1}, x_k] \\ i \neq0,n}} M_i\Delta x_i \leq
			\\
			2(n^* - 1) \cdot M\delta_1 + \suml_{k = 1}^n M_k^*\Delta x_k^* < \frac{\eps}{4} + \int_a^bf(x)dx + \frac{\eps}{4} = \int_a^bf(x)dx + \frac{\eps}{2}
		\end{multline*} 
	
	В случае без ограничений на $f$ верно, что $\exists m : f_1(x) = m + f(x) \geq 0$. Для этой функции справедливо соотношение
	\[
		U(P, f_1) < \int_a^bf_1(x)dx + \frac{\eps}{2}
	\]
	При этом есть 2 равенства:
	\[
		\System{
			&{U(P, f_1) = U(P, f) + m(b - a)}
			\\
			&{\int_a^b f_1(x)dx = \int_a^b f(x)dx + m(b - a)}
		}
	\]
	Отсюда получается требуемое утверждение. Теперь покажем, что
		\[
			(\forall \eps > 0)(\ \exists \delta_2 > 0) : (\forall P, \Delta(P) < \delta_2)\hookrightarrow\ \ L(P, f) >  \int_a^b f(x)dx - \frac{\eps}{2}
		\]
		Для этого положим $f_2 := -f$. Для которой верно по уже доказанному:
		\[
			(\exists \delta_2 > 0)\ (\forall P, \Delta P < \delta_2)\hookrightarrow U(P, f_2) < \int_a^bf_2(x)dx + \frac{\eps}{2} \Ra \int_a^bf(x)dx - \frac{\eps}{2} < L(P, f)
		\]
		Собeрем все в кучку:
		\[
			(\forall P : \Delta P < \delta = \min(\delta_1, \delta_2))\hookrightarrow U(P, f) - L(P, f) < \eps
		\]
		Так как любая интегральная сумма зажата между $L(P, f)$ и $U(P, f)$, то
		\[
			L(P, f) \leq  S(P, f, \{t_i\}) \leq U(P, f) \Ra \int_a^bf(x)dx = \liml_{\Delta(P) \to 0}S(P, f, \{t_i\})
		\]
	\end{itemize}
\end{proof}

\begin{note}
	Достаточность в доказанной теореме справедлива для интеграла Римана-Стилтьеса при дополнительном предположении, что $\alpha$ - непрерывна.
\end{note}

\begin{note}
	Если интеграл Римана-Стилтьеса определить как предел интегральных сумм, то свойство аддитивности не выполняется. Именно поэтому Лукашов предпочитает давать его через суммы Дарбу.
\end{note}


\begin{theorem}
	Если $\forall b' \in (a; b)\ \ f \in R[a; b']$ и $f$ ограничена на $[a; b]$, то $f \in R[a; b]$.
\end{theorem}

\begin{proof}
	Пусть $M := \sup\limits_{x \in [a; b]} |f(x)|$. Зафиксируем $\forall \eps > 0$. В качестве $b'$ положим $b' := b - \frac{\eps}{4M}$. Коль скоро функция интегрируема по Риману, то
	\[
		(\forall \eps > 0)\ (\exists P' \colon a = x_0 < \ldots < x_{n - 1} = b') :\ U(P', f) - L(P', f) < \frac{\eps}{2}
	\]
	Положим $P := P' \cup \{b\}$. Тогда
	\[
		U(P, f) - L(P, f) = U(P', f) - L(P', f) + \left(\sup\limits_{x \in [b'; b]} f(x) - \inf\limits_{x \in [b'; b]} f(x)\right) \frac{\eps}{4M} < \frac{\eps}{2} +2M\frac{\eps }{4M}=\eps
	\]
	Стало быть, $f \in R[a; b]$.
\end{proof}

\begin{corollary}
	Если $f$ непрерывна на $[a; b]$, кроме конечного числа точек, являющихся точками разрыва первого рода, то $f \in R[a; b]$.
\end{corollary}

\begin{definition}
	\textit{Интеграл с переменным верхним пределом} для $f \in R[a; b]$ это $F(x) = \int_a^{x} f(t)dt$ .
	
	Будем считать, что $F(a) = 0$, а для $\alpha > \beta\ \ \int_\alpha^\beta f(t)dt = -\int_\beta^\alpha f(t)dt$. 
\end{definition}

\begin{theorem} (Основные свойства интеграла с переменным верхним пределом)
	\begin{itemize}
		\item Если $f \in R[a; b]$, то $F(x) = \int_a^{x} f(t)dt$ непрерывна на $[a; b]$.
		
		\item Если, кроме того, $f$ непрерывна в точке $x_0 \in [a; b]$, то $F(x)$ дифференцируема в $x_0$, причём $F'(x_0) = f(x_0)$ (для $x_0 \in \{a, b\}$ нужно рассматривать соответствующую правую или левую производную).
	\end{itemize}
\end{theorem}

\begin{proof}~
	\begin{itemize}
		\item По свойству аддитивности, с учётом введённого понимания интеграла от большей абсциссы к меньшей:
		\[
			\forall x_1, x_2 \in [a; b],\ x_1 < x_2\ \ F(x_2) - F(x_1) = \int_a^{x_2} f(t)dt - \int_a^{x_1} f(t)dt= \int_{x_1}^{x_2} f(x)dx
		\]
		Оценим модуль этой разности (при условии, что $|f(t)| \le M$ и $0 < x_2 - x_1 < \frac{\eps}{M} = \delta$):
		\[
			|F(x_2) - F(x_1)| \le \int_{x_1}^{x_2} |f(x)|dx \le M(x_2 - x_1) < \eps
		\]
		Получается функция $F(x)$ равномерно непрерывна, а значит и просто непрерывна.
		
		\item Пусть теперь $f(x)$ непрерывна в $x_0\in[a,b]$. Посмотрим на разность следующего вида:
		\[
			\left|\frac{F(x) - F(x_0)}{x - x_0} - f(x_0)\right| = \left|\frac{1}{x - x_0} \int_{x_0}^x f(t)dt - f(x_0)\right| = \frac{1}{|x - x_0|} \cdot \left|\int_{x_0}^x (f(t) - f(x_0))dt\right|
		\]
		По определению непрерывности в точке $x_0$:
		\[
			(\forall \eps > 0)\ (\exists \delta > 0)\ :\ (\forall t, |t - x_0| < \delta,\ t \in [a;b]) \hookrightarrow |f(t) - f(x_0)| < \eps
		\]
		Откуда в свою очередь:
		\[
			\forall\ x : (0 < |x - x_0| < \delta, x \in [a; b]) \hookrightarrow \frac{1}{|x - x_0|} \cdot \left|\int_{x_0}^x (f(t) - f(x_0))dt\right| < \frac{1}{|x - x_0|} \cdot \left|\int_{x_0}^x \eps dx\right|  = \eps
		\]
		Мы получили, что 
		\begin{multline*}
			(\forall \eps > 0)\ (\exists \delta > 0)\ :\ (\forall x, 0 < |x - x_0| < \delta, x \in [a; b]) \hookrightarrow \left|\frac{F(x) - F(x_0)}{x - x_0} - f(x_0)\right| < \eps \lra
			\\
			\exists \liml_{x \to x_0} \frac{F(x) - F(x_0)}{x - x_0} = F'(x_0) = f(x_0)
		\end{multline*}
	\end{itemize}
\end{proof}

\begin{reminder}
	Функция $F$ называется первообразной для $f$ на $[a,b]$, если $(\forall x\in[a,b])\hookrightarrow F'(x)=f(x)$
\end{reminder}

\begin{theorem} (Основная теорема интегрального исчисления)
	Если $f \in R[a; b]$ имеет первообразную $F$ на $[a; b]$, то
	\[
		\int_a^b f(x)dx = F(b) - F(a) =: F(x)\Big|_a^b
	\]
	Это выражение также называется \textit{формулой Ньютона-Лейбница}.
\end{theorem}

\begin{proof}
	Зафиксируем разбиение $P \colon a = x_0 < x_1 < \ldots < x_n = b$. Распишем разность $F$:
	\[
		F(b) - F(a) = \suml_{k = 1}^n (F(x_k) - F(x_{k - 1}))
	\]
	По теореме Лагранжа $\exists t_k \in [x_{k - 1};x_k]\ F(x_k) - F(x_{k - 1}) = F'(t_k)\Delta x_k$. Тогда сумму можно переписать как
	\[
		\suml_{k = 1}^nF'(t_k)\Delta x_k = \suml_{k = 1}^nf(t_k)\Delta x_k = S(P, f, \{t_k\}) \xrightarrow[\Delta P \to 0]{} \int_a^b f(x)dx
	\]
	При этом $F(b) - F(a)$ от $P$ не зависит. Следовательно, теорема доказана.
\end{proof}

\begin{note}
	Формула Ньютона-Лейбница иногда принимается за определение интеграла. Который зачастую так и называется, интеграл Ньютона-Лейбница.
\end{note}

\begin{note}
	Если функция имеет первообразную, то она не обязана быть интегрируемой. Пример:
	\[
		F(x) = \System{
			&{x^2 \sin \frac{1}{x^2},\ x \neq 0}
			\\
			&{0,\ x = 0}
		}
	\]
	Отсюда $f(x) = F'(x)$:
	\[
		f(x) = \System{
			&{2x \sin \frac{1}{x^2} - \frac{2}{x}\cos \frac{1}{x^2},\ x \neq 0}
			\\
			&{0,\ x = 0}
		}
	\]
	Несложно увидеть, что $f(x)$ не может быть интегрируема по Риману на любом отрезке, содержащем 0, коль скоро тогда она не будет ограниченной.
\end{note}

\begin{note}
	Пример неинтегрируемой функции, имеющей первообразную, и при том ограниченной также есть. Но он гораздо хитрее и построить его сейчас нам не хватит средств. Т.к. он принадлежит Лебегу и имеет много идей из теории меры и его интеграла.
\end{note}